% !PRINTED 473
\weberpage{473}
Since, if we reserve to ourselves the determination of $x$, equation (3) can represent any arbitrary straight line, we can choose $x$ in such a way that among the four points $y, y', y'', y'''$ no two connected points occur. If we then set $u_y$ for the indeterminate quantity $t$ in (6), we obtain a rational equation

\weq{7}{\Psi (u_y, u_x, g, h) = 0,}

and this equation is no longer satisfied if for $(x)$ and $(y)$ we make a substitution from $G_{168}$ by which $u_x$ is changed and consequently $u_y$ is carried over into a value different from $u_y, u_{y'}, u_{y''}, u_{y'''}$.

Equation (7), regarded as an equation for $u_x$, therefore has only one root in common with the general resolvent $R_x$, and if one seeks the greatest common divisor of $R$ and $\Psi$, one obtains $u_x$ expressed rationally in terms of $u_y, g, h$. With this the general resolvent is reduced to the special one.

\begin{center}\rule{0.3\linewidth}{0.4pt}\end{center}

The group $G_{168}$ contains, further, a divisor of degree 21, $G_{21}$, which is generated by the substitutions $\chi, \tau$, and which gives rise to a resolvent of degree 8. As a root of this resolvent one can simply take the product $x_1 x_2 x_3$. The coefficients of this resolvent are invariants which rise up to degree 24. We shall not enter further into this resolvent here.

\begin{center}
\S.\,127.

\gesperrt{Permutation group of seven figures of degree 168.}
\end{center}

Since, as we have seen, the form problem of the ternary substitution group of degree 168 has a resolvent of degree 7, and the Galois resolvent of this equation of degree 7 is consequently of degree 168, it follows from this that in the general permutation group of seven figures, whose degree is $1.2.3.4.5.6.7 = 5040 = 168.30$, a divisor of degree 168 must be contained. The existence of this divisor was first
%CONT>>
